(1) 置\(x=\mu B/(k_BT)\)，单自旋配分函数\(z=e^x+e^{-x}=2\cosh x\)，\(Z=z^N\)。内能\(U=-\partial_\beta\ln Z=-N\mu B\tanh x\)，总磁化\(M=-\partial U_{\rm free}/\partial B=N\mu\tanh x\)（其中自由能\(F=-k_BT\ln Z\)）。在\(x\gg1\)时磁矩近乎全沿外场，在\(x\ll1\)时\(M\simeq N\mu^2B/(k_BT)\)恢复Curie响应。
(2) 熵\(S=-\partial F/\partial T=Nk_B[\ln(2\cosh x)-x\tanh x]\)，定场热容\(C_B=(\partial U/\partial T)_B=Nk_Bx^2\operatorname{sech}^2x\)。可用\(C_B=Nk_B(\beta\mu B)^2[1-\tanh^2x]\)独立核对，表明热容源于磁矩上下能级间的能量涨落而非连续动能。
(3) 对\(x>0\)，\(dS/dx=-Nk_Bx\operatorname{sech}^2x<0\)，故固定熵严格意味着\(x\)固定。由\(x=\mu B/(k_BT)\)得到\(B/T=\mathrm{const}\)，即理想顺磁盐可逆退磁时\(T_f=T_iB_f/B_i\)。磁化随退磁过程保持相同\(x\)，因此无需假设磁化恒零。
(4) 形式上令\(B_f\to0\)给\(T_f\to0\)，但在\(B=0\)有限\(T\)下\(z=2\)，残余熵\(Nk_B\ln2\)反映忽略相互作用的简并态。现实磁偶极、交换相互作用和核自旋效应会在极低温打破独立自旋模型，晶格又有正热容，所以不能凭理想比例律宣称有限步骤可达绝对零度。
